% Soft landing and inflation scares — academic seminar slides
% Replication / prototype talk for Bullard, Grimaud, Salle & Vermandel
%   Tinbergen Institute Discussion Paper TI 2025-001/VI (7 Jan 2025)
%   JME 157 (2026), 103871
%
% Compile from this directory:
%   pdflatex henk_seminar.tex
%   latexmk -pdf henk_seminar.tex
%
% Figure: figures/irf_costpush_fire_vs_sl.png (Python prototype IRF).
% Fallback search path: ../output/ if the PNG is regenerated there.
\documentclass[11pt,aspectratio=169]{beamer}

\usetheme{Madrid}
\usecolortheme{default}
\setbeamertemplate{navigation symbols}{}
\setbeamertemplate{itemize items}[circle]
\setbeamerfont{footnote}{size=\tiny}

\usepackage[T1]{fontenc}
\usepackage[utf8]{inputenc}
\usepackage{amsmath,amssymb}
\usepackage{booktabs}
\usepackage{graphicx}

\graphicspath{{figures/}{../output/}}

\newcommand{\paper}{Bullard, Grimaud, Salle and Vermandel (2025)}
\newcommand{\Esl}{E_t^{\mathrm{SL}}}
\newcommand{\Ere}{E_t}

\title[Soft landing \& inflation scares]{Soft Landing and Inflation Scares}
\subtitle{\texorpdfstring{Heterogeneous-expectation New Keynesian (HENK) dynamics\\
{\normalsize Replication / prototype seminar slides for \paper}}{HENK replication / prototype seminar slides}}
\author[Prototype slides]{\texorpdfstring{Open-source notes on the Tinbergen DP / JME model\\
{\small not the authors' replication package}}{Open-source prototype slides (not the authors' package)}}
\institute[MacroModellerBot]{\texttt{bullard-soft-landing/} in MacroModellerBot}
\date{September 2026}

\begin{document}

%----------------------------------------------------------------------
\begin{frame}[plain]
  \titlepage
  \vspace{-1.2em}
  \begin{center}
    {\footnotesize
    Source paper: \paper, Tinbergen Institute DP 2025-001\\
    (``Soft Landing and Inflation Scares''); JME 157 (2026), 103871.}
  \end{center}
\end{frame}

%----------------------------------------------------------------------
\begin{frame}{Outline}
  \tableofcontents
\end{frame}

%======================================================================
\section{Literature}
%----------------------------------------------------------------------
\begin{frame}{Why a soft landing?}
  \begin{block}{Core question \paper}
    Why did the post-pandemic U.S.\ inflation surge end in a
    \emph{soft landing}---disinflation without a major recession---unlike
    Volcker-era disinflation?
  \end{block}
  \begin{itemize}
    \item Standard FIRE New Keynesian models struggle to generate
          \emph{unanchored} long-run inflation beliefs after a transitory
          cost-push surge.
    \item The paper asks whether the \textbf{timing} and \textbf{strength}
          of the Fed's reaction matter when long-run inflation expectations
          are \textbf{heterogeneous} and can lose their anchor to the target.
    \item That unanchoring is an \textbf{inflation scare} in the sense of
          Goodfriend (1993): a rise in perceived trend inflation that
          raises the real cost of subsequent disinflation.
  \end{itemize}
\end{frame}

%----------------------------------------------------------------------
\begin{frame}{Soft landing vs.\ inflation-scare literature}
  \begin{itemize}
    \item \textbf{Goodfriend (1993).}
          Inflation scares are shifts in long-run inflation beliefs that
          force the central bank to choose between validating the scare
          and inflicting a recession. The 1980s disinflations are the
          classic case.
    \item \textbf{Post-COVID debate.}
          Was 2021--23 a transitory relative-price shock (FIRE-like), or
          a scare in which the distribution of trend-inflation beliefs
          moved? \paper{} take the second view.
    \item \textbf{Policy timing.}
          The abstract's claim is that \emph{when} the Fed tightened
          matters more than how large a later hike would have been for
          containing the scare (WP \S4 counterfactuals).
    \item \textbf{This talk.}
          Situates that argument in a 3-equation HENK block, then reports
          what the open-source Python/Dynare prototype in this repo
          actually implements---not the authors' estimated Table~3 paths.
  \end{itemize}
\end{frame}

%----------------------------------------------------------------------
\begin{frame}{FIRE vs.\ learning / heterogeneous expectations}
  \begin{itemize}
    \item \textbf{FIRE NK.}
          A single rational forecast; $\phi_t\equiv 0$. Cost-push inflation
          is temporary once the shock decays. Nested as the paper's
          homogeneous-belief special case.
    \item \textbf{Adaptive / statistical learning}
          (Evans and Honkapohja 2001; Hommes 2021 survey):
          agents estimate perceived laws of motion; mean beliefs can drift.
    \item \textbf{Social learning (SL)}
          (Grimaud, Salle and Vermandel 2025, JEDC toolbox paper):
          a finite population copies better-performing inflation intercepts
          (news $\to$ fitness $\to$ tournament). Dispersion is an
          \emph{object}, not a residual.
    \item \textbf{HENK contribution.}
          Only long-run \emph{inflation} beliefs are non-FIRE; output-gap
          expectations in the IS curve remain model-consistent
          (internal rationality / quasi-RE observer).
    \item Closed form retains the \emph{entire} time-varying cross-section
          of subjective intercepts $\{a_{j,t}\}$, not just the mean.
  \end{itemize}
\end{frame}

%----------------------------------------------------------------------
\begin{frame}{Related companions}
  \begin{itemize}
    \item \textbf{Toolbox.}
          Grimaud, Salle and Vermandel (2025), \emph{A Dynare toolbox for
          social learning expectations}, JEDC; Mendeley
          \texttt{fzmx3vkt66}. Authors solve HENK with this stack.
    \item \textbf{Communication / ELB.}
          Arifovic, Grimaud, Salle and Vermandel (\emph{JMCB}):
          central-bank announcements in the SL fitness stage as an
          anchoring device---a natural companion to rate-hike timing.
    \item \textbf{Contrast, not synthesis} (NOTES.md digest):
          Kaplan and Miyahara (NBER WP 35400)---HANK plus state-dependent
          pricing; Angeletos, Lian, Wolf and Zhang (NBER WP 35642)---
          fiscal-monetary interactions under finite-horizon refinements.
          Those papers stress \emph{balance-sheet} heterogeneity; HENK
          stresses \emph{belief} heterogeneity in an otherwise RANK NK.
    \item Experience-based anchoring (Malmendier--Nagel style) and
          finite planning horizons are \emph{not} in the estimated HENK
          block; they sit in the same unanchoring queue (see extensions).
  \end{itemize}
\end{frame}

%======================================================================
\section{The model}
%----------------------------------------------------------------------
\begin{frame}{HENK at a glance}
  \begin{itemize}
    \item Stylized NK economy: NKPC + Euler/IS + Taylor rule with smoothing.
    \item Population of $J$ agents (even), identical except for idiosyncratic
          long-run inflation intercepts $a_{j,t}$.
    \item Aggregate SL inflation forecast $=$ cross-sectional average of
          individual forecasts.
    \item \textbf{Only departure from RE:} information friction on long-run
          inflation via social learning. Output-gap leads stay rational.
    \item Hats $=$ deviations from steady state / inflation target.
  \end{itemize}
  \vfill
  \begin{block}{Belief wedge (WP eq.\ 8)}
    \[
      \Esl(\hat\pi_{t+1})
      = \phi_t + \Ere(\hat\pi_{t+1}),
      \qquad
      \phi_t \equiv \tfrac1J\sum_{j=1}^J a_{j,t}.
    \]
    Nested FIRE: $\phi_t=0$ for all $t$ $\Rightarrow$ $\Esl=\Ere$.
  \end{block}
\end{frame}

%----------------------------------------------------------------------
\begin{frame}{Aggregate block (WP eqs.\ 1--3)}
  \begin{align}
    \hat\pi_t
      &= \kappa\,\hat y_t
         + \beta\,\Esl(\hat\pi_{t+1})
         + u_t
      \tag{NKPC} \\[0.4em]
    \hat y_t
      &= \Ere(\hat y_{t+1})
         - \sigma^{-1}\bigl(
             \hat\iota_t - \Esl(\hat\pi_{t+1})
           \bigr)
         + g_t
      \tag{IS} \\[0.4em]
    \hat\iota_t
      &= \rho_\iota\,\hat\iota_{t-1}
         + (1-\rho_\iota)
           \bigl(\phi_\pi\,\hat\pi_t + \phi_y\,\hat y_t\bigr)
         + v_t
      \tag{Taylor}
  \end{align}
  \begin{itemize}
    \item Cost-push $u_t$ and demand $g_t$ are ARMA(1,1); monetary $v_t$
          is white noise (WP \S2.4).
    \item Seminar-level reading: a scare ($\phi_t\uparrow$) acts like a
          persistent inflation shock in the NKPC \emph{and} eases the
          ex-ante real rate in the IS curve.
  \end{itemize}
\end{frame}

%----------------------------------------------------------------------
\begin{frame}{Martingale beliefs and the PLM}
  \begin{block}{Assumption 2 (WP)}
    \[
      \Ere(\phi_{t+1})=\phi_t
      \qquad\text{(random-walk / martingale mean belief).}
    \]
  \end{block}
  \begin{itemize}
    \item Individual perceived law of motion (WP eq.\ 6):
          \[
            \hat\pi_t^{(j)}
            = a_{j,t}
              + P_{1,\bullet}\,\hat z_{t-1}
              + \widetilde Q_{1,\bullet}\,E_t.
          \]
    \item Full MSV rows $P,\widetilde Q$ come from the Dynare SL toolbox
          (\texttt{SL\_get\_pq}); they are \textbf{not} recovered from the
          WP text alone.
    \item Prototype stand-in for fitness (documented provisional):
          \[
            \mathbb{E}[\pi_s\mid m]
            = m + A_\pi i_{s-1} + D_\pi u_s + G_\pi g_s + B_\pi v_s.
          \]
  \end{itemize}
\end{frame}

%----------------------------------------------------------------------
\begin{frame}{Social learning (WP eqs.\ 10--12)}
  Three-step recursive process each period, $j=1,\ldots,J$:
  \begin{enumerate}
    \item \textbf{News / mutation} (Calvo-like probability $\mu$):
          \[
            m_{j,t}
            = a_{j,t-1}
              + \mathbf{1}_{\{\varpi_{j,t}\le\mu\}}
                \,(\iota_{j,t}+\lambda_t).
          \]
    \item \textbf{Fitness:} discounted squared forecast errors on inflation
          history (decay $\rho$).
    \item \textbf{Tournament:} agents paired; both copy the more accurate
          intercept.
  \end{enumerate}
  \vspace{0.4em}
  Aggregate update $\phi_t=\phi_{t-1}+S(\cdot)$ with $S$ nonlinear
  (no closed form for $S$). Authors solve it with the GSV Dynare SL
  toolbox; this repo's \texttt{social\_learning.py} implements the
  three micro steps on a finite $J$.
\end{frame}

%----------------------------------------------------------------------
\begin{frame}{Shocks and estimation (paper)}
  \begin{itemize}
    \item Demand: $g_t=\rho_g g_{t-1}+\varepsilon_t^g-\mu_g\varepsilon_{t-1}^g$.
    \item Cost-push: $u_t=\rho_u u_{t-1}+\varepsilon_t^u-\mu_u\varepsilon_{t-1}^u$.
    \item Monetary: $v_t=\varepsilon_t^v$.
    \item SL news: aggregate $\lambda_t$, idiosyncratic $\iota_{j,t}$.
    \item Bayesian estimation with an inversion filter
          (Cuba-Borda, Guerrieri, Iacoviello and Zhong 2019) on
          \textbf{1985Q1--2023Q4} U.S.\ macros \emph{and} SPF inflation
          forecasts (WP \S3). $\beta=0.99$ calibrated.
  \end{itemize}
  \vspace{0.3em}
  {\footnotesize
  \begin{tabular}{@{}lrlr@{}}
    \toprule
    Param. & Table 1 mean & Param. & Table 1 mean \\
    \midrule
    $\kappa$ & 0.2032 & $\rho_\iota$ & 0.8179 \\
    $\sigma$ & 1.6993 & $\rho_u,\mu_u$ & 0.6328, 0.4887 \\
    $\phi_\pi$ & 1.7131 & $\mu$ (news) & 0.4357 \\
    $\phi_y$ & 0.1564 & $\rho$ (fitness) & 0.7745 \\
    \bottomrule
  \end{tabular}
  }\\[0.4em]
  {\scriptsize Use Table 1 posterior means (not the WP narrative aside
  $\phi_\pi\approx 1.82$). Estimation itself is \textbf{not} in this repo.}
\end{frame}

%======================================================================
\section{Contributions}
%----------------------------------------------------------------------
\begin{frame}{What the Tinbergen / JME paper adds}
  Relative to FIRE NK and to the inflation-scare literature:
  \begin{itemize}
    \item A \textbf{micro-founded SL} layer that can \emph{generate}
          inflation scares: heterogeneous long-run intercepts, time-varying
          dispersion, possible unanchoring from target.
    \item A \textbf{closed-form} representation that keeps the full
          cross-section of beliefs, then Bayesian estimation on macros+SPF.
    \item Quantitative claim: \textbf{timing} of tightening dominates
          delayed \textbf{strength} for containing the scare (conditional
          on policy not being too weak).
    \item Historical narrative (WP \S4): the Fed was behind the curve in
          2021; earlier tightening could have cut the inflation peak
          without a recession; further delay risks entrenchment and
          larger output losses.
    \item Soft landing as successful (if late) scare management under
          HENK---not as FIRE-only mean reversion.
  \end{itemize}
\end{frame}

%----------------------------------------------------------------------
\begin{frame}{What this open-source prototype adds}
  Honest split vs.\ the published paper:
  \begin{itemize}
    \item \textbf{Runnable Python} FIRE nesting, finite-$J$ SL micro block,
          coupled HENK IRFs, and illustrative timing-vs-strength
          counterfactuals---no Matlab required.
    \item \textbf{Dynare scaffold} for the linear NK+ARMA block, aimed at
          the GSV SL toolbox once \texttt{TOOLBOX\_PATH} is set.
    \item Transparent, documented approximations (quasi-RE observer,
          martingale $\phi$, provisional $C,D,G$ loadings).
    \item \textbf{Does not} reproduce author Table~3 historical-shock
          paths, inversion-filter estimation, or SPF measurement equations.
    \item No official replication package for the paper was found
          (GitHub / journal supplement / author download).
  \end{itemize}
  \vfill
  {\small This talk is a map of the \emph{paper} plus a status report on
  the \emph{repo}---two different objects.}
\end{frame}

%======================================================================
\section{Results and discussion}
%----------------------------------------------------------------------
\begin{frame}{Prototype IRF: cost-push, FIRE vs.\ SL}
  \begin{center}
    \includegraphics[width=\textwidth,height=0.68\textheight,keepaspectratio]%
      {irf_costpush_fire_vs_sl.png}
  \end{center}
  {\scriptsize Python HENK prototype; unit cost-push impulse; Table~1
  params; seed $=1$, $J=100$, ARMA(1,1) on; horizon 40.
  \textbf{[PROVISIONAL]} not the authors' Dynare-toolbox IRFs.}
\end{frame}

%----------------------------------------------------------------------
\begin{frame}{Reading the cost-push IRF}
  Qualitative pattern in \texttt{output/fire\_vs\_sl\_summary.txt}
  (and the figure):
  \begin{itemize}
    \item \textbf{FIRE} ($\phi\equiv 0$): inflation jumps on impact
          (unit shock: $\pi_0\approx 1.39$) then mean-reverts quickly---
          even slightly undershoots. Temporary cost-push, as in a
          standard 3-eq NK.
    \item \textbf{SL}: impact almost identical (beliefs have not had
          time to move). Thereafter $\phi_t$ drifts away from 0 and
          inflation stays \emph{higher for longer} (turns positive again
          while FIRE is still decaying).
    \item Mechanism: the wedge $\phi_t$ feeds the NKPC and the IS real
          rate; dispersed long-run beliefs amplify and persist the scare.
    \item Same structural parameters in both regimes---the extra
          persistence is \emph{not} a more persistent $u_t$.
  \end{itemize}
\end{frame}

%----------------------------------------------------------------------
\begin{frame}{Timing vs.\ strength (prototype only)}
  Shared 12-quarter cost-push surge ($0.01$ innovations), identical SL
  seed, delay $=8$ quarters
  (\texttt{timing\_counterfactual.py}; NOTES.md STATUS):
  \begin{center}
  {\footnotesize
  \begin{tabular}{@{}lccc@{}}
    \toprule
    Scenario & peak $\pi$ & min $y$ & $\max|\phi|$ \\
    \midrule
    Baseline (Table 1) & 0.093 & $-0.169$ & 0.037 \\
    Earlier ($\phi_\pi{+}10\%$, $\rho_\iota\times 0.9$ from $t=0$)
      & 0.068 & $-0.158$ & 0.034 \\
    Stronger delayed ($\phi_\pi{+}10\%$ from $t=8$)
      & 0.075 & $-0.171$ & 0.037 \\
    \bottomrule
  \end{tabular}
  }
  \end{center}
  \begin{itemize}
    \item Pattern is \emph{consistent with} WP \S4: earlier tightening
          beats delayed strength for peak inflation / scare containment.
    \item \textbf{Not} a claim of matching author Table~3 or estimated
          historical shocks. Peaks occur at the horizon end; with a
          well-specified PLM, $\phi$ paths can be news-dominated.
  \end{itemize}
\end{frame}

%----------------------------------------------------------------------
\begin{frame}{Why this matters for 2021--22}
  \begin{itemize}
    \item \textbf{Post-COVID surge.}
          A sequence of cost-push / relative-price shocks plus delayed
          liftoff is exactly the environment in which an inflation scare
          can develop if long-run beliefs are heterogeneous.
    \item \textbf{Anchoring is an object.}
          Under FIRE, ``transitory'' is almost automatic. Under SL,
          anchoring is a \emph{distribution} that policy timing can
          protect or lose.
    \item \textbf{Fed timing.}
          Paper: earlier 2021 tightening could have reduced the inflation
          peak without a recession; further delay would have unanchored
          expectations and raised output losses. Soft landing $\approx$
          containing the scare late rather than never.
    \item \textbf{Communication} (companion JMCB idea): talking the
          target into the tournament is a potential substitute for
          earlier hikes---not coded here.
  \end{itemize}
\end{frame}

%======================================================================
\section{What the replicated model achieves}
%----------------------------------------------------------------------
\begin{frame}{Runnable artefacts}
  \begin{itemize}
    \item \texttt{src/re\_nk\_irf.py} --- FIRE / RE 3-eq monetary IRFs;
          $\phi_t=0$; Table~1 params. Nested by the HENK simulator.
    \item \texttt{src/social\_learning.py} --- finite-$J$ news, fitness,
          tournament (WP eqs.\ 10--12).
    \item \texttt{src/henk\_sim.py} --- FIRE vs.\ SL IRFs with ARMA(1,1)
          $u,g$ and richer PLM common forecast
          $\to$ \texttt{output/irf\_*\_\{fire,sl,sl\_minus\_fire\}.csv}.
    \item \texttt{src/timing\_counterfactual.py} --- baseline / earlier /
          stronger-delayed under a shared surge
          $\to$ \texttt{output/counterfactuals/}.
    \item \texttt{src/plot\_costpush\_irf.py} --- figure on the previous
          IRF slide.
  \end{itemize}
  FIRE paths nest \texttt{re\_nk\_irf} on the monetary $A,B$ block.
\end{frame}

%----------------------------------------------------------------------
\begin{frame}{Calibration, Dynare scaffold, caveats}
  \begin{itemize}
    \item \textbf{Calibration:} Tinbergen Table~1 posterior means
          ($\kappa,\sigma,\phi_\pi,\phi_y,\rho_\iota$, ARMA, SL news/fitness)
          plus $\beta=0.99$. Shock stds live in the Dynare \texttt{.mod};
          Python IRFs default to \emph{unit} impulses.
    \item \textbf{Dynare:} \texttt{dynare/bullard\_henk.mod} is the
          FIRE/RE nesting of WP (1)--(3). Driver
          \texttt{run\_bullard\_henk.m} calls \texttt{SL\_get\_pq} then
          \texttt{SL\_chain} \emph{after} you set \texttt{TOOLBOX\_PATH}
          to the unzipped Mendeley \texttt{V3/Toolbox}.
    \item \textbf{Not yet aligned:} Python $C,D,G$ loadings vs.\ toolbox
          $P,\widetilde Q,R$; inflation-only SL via
          \texttt{mut\_sd}$_y=0$ is provisional; no separate $\lambda$
          shock in toolbox options.
    \item \textbf{Blocked:} inversion-filter estimation, SPF observables,
          author Table~3, official replication files.
  \end{itemize}
\end{frame}

%======================================================================
\section{Simplifying assumptions}
%----------------------------------------------------------------------
\begin{frame}{What is on, what is shut off}
  \textbf{In the paper's baseline HENK (and this prototype):}
  \begin{itemize}
    \item 3-equation NK block (no capital, no HANK, no open economy).
    \item Heterogeneity \emph{only} in long-run inflation intercepts;
          output-gap expectations remain RE / model-consistent.
    \item Steady-state learning around a zero-deviation SS; hats vs.\
          target.
    \item Quasi-RE observer: each period SL sets $\phi_t$ from \emph{past}
          inflation, then the linear macro block takes $\phi_t$ as given
          (martingale $\mathbb{E}\phi'=\phi$). Not a simultaneous
          nonlinear RE+SL fixed point.
  \end{itemize}
  \textbf{Shut off / not in the prototype:}
  \begin{itemize}
    \item Own-cost-belief NKPC extensions; Rotemberg microfoundations
          of $\kappa$ (App.~B); MA lag in MSV loadings (paths are ARMA,
          loadings still AR(1) expectations).
    \item Delayed-strength loadings switch without agents anticipating
          the future rule.
  \end{itemize}
\end{frame}

%======================================================================
\section{Extensions}
%----------------------------------------------------------------------
\begin{frame}{Future developments that fit this codebase}
  \begin{enumerate}
    \item \textbf{Dynare SL alignment.}
          Install GSV toolbox; set \texttt{TOOLBOX\_PATH}; cross-check
          Python $C,D,G$ against \texttt{mo.PP/QQ/RR}; optional MA state
          in loadings.
    \item \textbf{Phillips-block extensions.}
          Own-cost-belief NKPC; App.~B Rotemberg $\leftrightarrow\kappa$.
    \item \textbf{Estimation stack} (blocked on data/codes).
          Inversion filter, SPF measurement (WP \S3.1 / App.~A),
          historical-shock Table~3.
    \item \textbf{Companions in the literature queue.}
          Experience-based anchoring; finite planning horizons;
          CB communication in the tournament (Arifovic et al.);
          UK/BoE timing analogue (NOTES.md).
    \item \textbf{Later synthesis} (not this repo): HENK beliefs + HANK
          fiscal non-Ricardians for a fuller post-COVID accounting.
  \end{enumerate}
\end{frame}

%----------------------------------------------------------------------
\begin{frame}{Takeaways}
  \begin{itemize}
    \item Inflation scares are a \emph{distributional} object: dispersed
          long-run inflation beliefs can persist after the cost-push dies.
    \item FIRE NK is nested ($\phi=0$) and looks transitory; SL adds
          persistence through $\phi_t$.
    \item The paper's policy punchline is \textbf{timing over delayed
          strength}. The Python prototype is qualitatively consistent,
          not a numerical replica of Table~3.
    \item This repo is a transparent stepping stone: FIRE+SL IRFs,
          a Dynare scaffold, and documented gaps
          (toolbox, estimation, author codes).
  \end{itemize}
  \vfill
  {\scriptsize
  \textbf{Selected references.}
  Bullard, Grimaud, Salle and Vermandel (2025), Tinbergen DP 2025-001 /
  JME (2026);
  Goodfriend (1993);
  Grimaud, Salle and Vermandel (2025), JEDC;
  Evans and Honkapohja (2001);
  Cuba-Borda et al.\ (2019);
  Arifovic, Grimaud, Salle and Vermandel (\emph{JMCB}).
  Compile: \texttt{pdflatex henk\_seminar.tex} from
  \texttt{bullard-soft-landing/beamer/}.}
\end{frame}

\end{document}
